\documentclass[a4paper,11pt]{article}
\usepackage[utf8]{inputenc}
\usepackage[T1]{fontenc}
\usepackage[portuguese]{babel}
\usepackage{geometry}
\geometry{a4paper, margin=2.5cm}
\usepackage{hyperref}
\usepackage{listings}
\usepackage{xcolor}
\usepackage{graphicx}
\usepackage{setspace}
\onehalfspacing

% Hyperlink colors
\hypersetup{
    colorlinks=true,
    linkcolor=blue,
    urlcolor=blue,
    citecolor=blue
}

\title{Cyber-Toolbox \\ Documentação do Código-Fonte}
\author{Martinho Caeiro (23917) \\ Mestrado de Engenharia em Segurança Informática \\ Linguagens de Programação Dinâmicas}
\date{\today}

\begin{document}

\maketitle
\tableofcontents
\newpage

\section{Introdução}
Este documento contém a documentação técnica extraída do código-fonte do projeto Cyber-Toolbox.
Todos os scripts foram desenvolvidos em Python 3 com comentários em inglês e mensagens de saída em português (pt-PT).

\section{Menu Principal}

\subsection{menu\_launcher}
\label{sec:menu_launcher}

\textbf{Ficheiro:} \texttt{src/menu\_launcher.py}

\textbf{Autor:} Martinho Caeiro (23917)

\textbf{Descrição:}

Simple menu launcher to manage/run scripts in the ./scripts folder.

\textbf{Uso:}
\begin{verbatim}
python3 menu_launcher.py
\end{verbatim}

\textbf{Exemplo:}
\begin{verbatim}
python3 menu_launcher.py
\end{verbatim}

\textbf{Secções do código:}
\begin{itemize}
    \item Helpers
    \item Script execution
    \item Main loop
\end{itemize}

\textbf{Funções e Classes principais:}
\begin{itemize}
    \item \texttt{ensure\_dirs()}
    \item \texttt{clear\_screen()}
    \item \texttt{show\_menu()}
    \item \texttt{prompt\_args(arg\_defs)}
    \item \texttt{run\_script(path, args)}
    \item \texttt{main\_menu()}
\end{itemize}

\newpage
\subsection{scripts\_config}
\label{sec:scripts_config}

\textbf{Ficheiro:} \texttt{src/scripts\_config.py}

\textbf{Autor:} Martinho Caeiro (23917)

\textbf{Descrição:}

Defines the configuration of available scripts, including names, files,
and required arguments.

\newpage
\section{Scripts Individuais}

\subsection{log\_analyzer}
\label{sec:log_analyzer}

\textbf{Ficheiro:} \texttt{src/scripts/log\_analyzer.py}

\textbf{Autor:} Martinho Caeiro (23917)

\textbf{Descrição:}

Analyze Apache/Nginx access logs and OpenSSH auth logs.

\textbf{Funcionalidades:}
\begin{itemize}
    \item Parse Apache/Nginx combined log lines and OpenSSH auth.log lines.
    \item Extract IP, timestamp, service, status (success/fail), and details.
    \item Optionally resolve country via GeoLite2 DB (pass --geoip-db).
    \item Emit JSON and CSV reports grouped per event.
\end{itemize}

\textbf{Uso:}
\begin{verbatim}
python log_analyzer.py --files access.log auth.log --outdir reports --geoip-db /path/GeoLite2-Country.mmdb
\end{verbatim}

\textbf{Exemplo:}
\begin{verbatim}
python log_analyzer.py --files access.log auth.log --outdir reports
\end{verbatim}

\textbf{Secções do código:}
\begin{itemize}
    \item Dependencies
    \item Patterns
    \item Resolvers
    \item Parsers
    \item Reporting
    \item CLI
\end{itemize}

\textbf{Funções e Classes principais:}
\begin{itemize}
    \item Classe: \texttt{GeoIPResolver}
    \item \texttt{parse\_apache\_line(line)}
    \item \texttt{parse\_ssh\_line(line)}
    \item \texttt{parse\_ufw\_line(line)}
    \item \texttt{analyze\_files(file\_paths, geoip\_db, outdir)}
    \item \texttt{find\_default\_geoip\_db()}
    \\ Look for GeoLite2 DB files inside the repository `src/data` folder.
    \item \texttt{summarize\_events(events)}
    \\ Print a small summary: counts per service and top countries.
    \item \texttt{main(argv)}
    \item \texttt{country\_for(self, ip)}
\end{itemize}

\newpage
\subsection{messenger}
\label{sec:messenger}

\textbf{Ficheiro:} \texttt{src/scripts/messenger.py}

\textbf{Autor:} Martinho Caeiro (23917)

\textbf{Descrição:}

Simple secure messenger (server + client CLI) for the Cyber-Toolbox project.

\textbf{Funcionalidades:}
\begin{itemize}
    \item Server generates an RSA keypair (data/keys/) with init-keys.
    \item Clients connect to the server (TCP) and send JSON commands.
    \item Messages are stored on the server only, in data/messages/.
    \item Each message is hybrid-encrypted: a Fernet payload + the Fernet key
    \item encrypted with the server RSA public key.
    \item Clients can list/download/delete messages where they are participants.
    \item Server can produce backups (symmetric password-protected or asymmetric).
\end{itemize}

\textbf{Uso:}
\begin{verbatim}
python messenger.py init-keys
python messenger.py serve --host 127.0.0.1 --port 9009
python messenger.py send --from alice --to bob --message "Olá"
\end{verbatim}

\textbf{Exemplo:}
\begin{verbatim}
python messenger.py list --user alice
\end{verbatim}

\textbf{Nota:}

This script is intentionally simple and synchronous. It is meant for local
testing and demonstration, not production use.

\textbf{Secções do código:}
\begin{itemize}
    \item Dependencies
    \item Key management
    \item Message storage
    \item Server
    \item Client \& CLI
\end{itemize}

\textbf{Funções e Classes principais:}
\begin{itemize}
    \item \texttt{ensure\_dirs()}
    \item \texttt{generate\_rsa\_keys(key\_size)}
    \item \texttt{load\_private\_key()}
    \item \texttt{load\_public\_key()}
    \item \texttt{hybrid\_encrypt(plaintext)}
    \\ Encrypt plaintext with a fresh Fernet key, then encrypt that key with RSA public key.
    \item \texttt{hybrid\_decrypt(b64\_enc\_sym\_key, b64\_payload)}
    \item \texttt{store\_message(sender, recipients, text)}
    \item \texttt{list\_messages\_for(user)}
    \item \texttt{load\_message\_file(mid)}
    \item \texttt{delete\_message(mid, actor)}
    \item \texttt{backup\_all\_sym(password)}
    \\ Create a password-protected backup file (Fernet with key derived from password).
    \item \texttt{backup\_all\_asym()}
    \\ Create an asymmetric backup encrypted with server public key.
    \item Classe: \texttt{ThreadedTCPHandler}
    \item \texttt{run\_server(host, port)}
    \item \texttt{client\_send(host, port, payload)}
    \item \texttt{print\_response(resp)}
    \\ Print a response payload in pt-PT.
    \item \texttt{main()}
    \item \texttt{handle(self)}
\end{itemize}

\newpage
\subsection{password\_manager}
\label{sec:password_manager}

\textbf{Ficheiro:} \texttt{src/scripts/password\_manager.py}

\textbf{Autor:} Martinho Caeiro (23917)

\textbf{Descrição:}

Simple password manager with RSA + Fernet encryption and TOTP 2FA.

\textbf{Uso:}
\begin{verbatim}
python3 password_manager.py --init
python3 password_manager.py --create --url example.com --user alice --pass s3cr3t
python3 password_manager.py --list --otp 123456
python3 password_manager.py --view <id> --otp 123456
python3 password_manager.py --update <id> [--url ..] [--user ..] [--pass ..] --otp 123456
python3 password_manager.py --delete <id> --otp 123456
\end{verbatim}

\textbf{Exemplo:}
\begin{verbatim}
python3 password_manager.py --list --otp 123456
Files created (in the same folder):
- private_key.pem  (PEM, kept secret)
- public_key.pem   (PEM)
- totp_secret.txt  (base32 secret for TOTP)
- records.json     (encrypted records database)
\end{verbatim}

\textbf{Nota:}

This is a simple educational implementation. Do not use for production secrets.

\textbf{Secções do código:}
\begin{itemize}
    \item Dependencies
    \item Paths
    \item Key management
    \item Commands
    \item CLI
\end{itemize}

\textbf{Funções e Classes principais:}
\begin{itemize}
    \item \texttt{ensure\_files\_dir()}
    \item \texttt{init\_keys\_and\_totp()}
    \\ Generate RSA keypair and a new TOTP secret.
    \item \texttt{load\_public\_key()}
    \item \texttt{load\_private\_key()}
    \item \texttt{load\_records()}
    \item \texttt{save\_records(recs)}
    \item \texttt{encrypt\_record\_plaintext(plaintext\_bytes)}
    \\ Encrypt plaintext bytes using Fernet and encrypt the key with RSA public key.
    \item \texttt{decrypt\_record(enc\_cipher\_b64, enc\_key\_b64)}
    \item \texttt{require\_otp(code)}
    \item \texttt{cmd\_create(args)}
    \item \texttt{cmd\_list(args)}
    \item \texttt{cmd\_view(args)}
    \item \texttt{cmd\_delete(args)}
    \item \texttt{cmd\_update(args)}
    \item \texttt{parse\_args()}
    \item \texttt{main()}
    \item \texttt{run\_interactive\_menu()}
    \item Classe: \texttt{A}
    \item Classe: \texttt{A}
    \item Classe: \texttt{A}
    \item Classe: \texttt{A}
    \item Classe: \texttt{A}
\end{itemize}

\newpage
\subsection{port\_knocker}
\label{sec:port_knocker}

\textbf{Ficheiro:} \texttt{src/scripts/port\_knocker.py}

\textbf{Autor:} Martinho Caeiro (23917)

\textbf{Descrição:}

Sends a sequence of packets to specific ports of a remote server to enable SSH access.

\textbf{Uso:}
\begin{verbatim}
python3 port_knocker.py <host> <port1> <port2> <port3>
\end{verbatim}

\textbf{Exemplo:}
\begin{verbatim}
python3 port_knocker.py 192.168.1.100 7000 8000 9000
\end{verbatim}

\textbf{Secções do código:}
\begin{itemize}
    \item Port knocking
    \item Main loop
\end{itemize}

\textbf{Funções e Classes principais:}
\begin{itemize}
    \item \texttt{send\_knock(host, port)}
    \item \texttt{main()}
\end{itemize}

\newpage
\subsection{port\_scanner}
\label{sec:port_scanner}

\textbf{Ficheiro:} \texttt{src/scripts/port\_scanner.py}

\textbf{Autor:} Martinho Caeiro (23917)

\textbf{Descrição:}

Scans ports on multiple targets to identify which are open.

\textbf{Uso:}
\begin{verbatim}
python3 port_scanner.py <target1,target2,...> <start_port> <end_port>
\end{verbatim}

\textbf{Exemplo:}
\begin{verbatim}
python3 port_scanner.py target1.com,target2.com,127.0.0.1 1 1024
\end{verbatim}

\textbf{Secções do código:}
\begin{itemize}
    \item Scanner
    \item CLI
\end{itemize}

\textbf{Funções e Classes principais:}
\begin{itemize}
    \item \texttt{scan\_host(host, start\_port, end\_port, timeout)}
    \\ Scan ports on a specific host.
    \item \texttt{main()}
\end{itemize}

\newpage
\subsection{syn\_flooder}
\label{sec:syn_flooder}

\textbf{Ficheiro:} \texttt{src/scripts/syn\_flooder.py}

\textbf{Autor:} Martinho Caeiro (23917)

\textbf{Descrição:}

Sends TCP SYN packets to simulate a denial-of-service attack (SYN Flood/DoS).
Uses Scapy for low-level packet crafting.

\textbf{Uso:}
\begin{verbatim}
python3 syn_flooder.py [--target HOST] [--port PORT] [--duration SECONDS] [--threads N]
\end{verbatim}

\textbf{Exemplo:}
\begin{verbatim}
python3 syn_flooder.py --target 192.168.1.1 --port 80 --duration 10
\end{verbatim}

\textbf{Secções do código:}
\begin{itemize}
    \item SYN Flooder functions
\end{itemize}

\textbf{Funções e Classes principais:}
\begin{itemize}
    \item \texttt{generate\_random\_ip()}
    \\ Generate a random spoofed source IP address.
    \item \texttt{generate\_random\_port()}
    \\ Generate a random source port.
    \item \texttt{syn\_flood\_worker(target, target\_port, duration, packets\_sent, lock)}
    \\ Worker thread that sends TCP SYN packets to a target.
    \item \texttt{flood\_target(target, target\_port, duration, num\_threads)}
    \\ Send SYN flood with multiple threads.
    \item \texttt{main()}
\end{itemize}

\newpage
\subsection{udp\_flooder}
\label{sec:udp_flooder}

\textbf{Ficheiro:} \texttt{src/scripts/udp\_flooder.py}

\textbf{Autor:} Martinho Caeiro (23917)

\textbf{Descrição:}

Sends UDP packets to a target (UDP Flood attack).

\textbf{Uso:}
\begin{verbatim}
python3 udp_flooder.py [--target HOST] [--port PORT] [--duration SECONDS] [--threads N]
\end{verbatim}

\textbf{Exemplo:}
\begin{verbatim}
python3 udp_flooder.py --target 192.168.1.1 --port 53 --duration 10
\end{verbatim}

\textbf{Secções do código:}
\begin{itemize}
    \item Workers
    \item Flood logic
    \item CLI
\end{itemize}

\textbf{Funções e Classes principais:}
\begin{itemize}
    \item \texttt{udp\_flood\_worker(target, port, packet\_size, duration, packets\_sent)}
    \\ Worker thread that sends UDP packets to a target.
    \item \texttt{flood\_target(target, port, packet\_size, duration, num\_threads)}
    \\ Send UDP flood with multiple threads.
    \item \texttt{main()}
\end{itemize}

\newpage
\end{document}